% SPDX-FileCopyrightText: Copyright © 2026 Idiap Research Institute <contact@idiap.ch>
% SPDX-FileContributor: Andre Freitas <andre.freitas@idiap.ch>
% SPDX-License-Identifier: MIT
% Neuro-symbolic AI Group
% File Documentation
%
% Motivation:
%   Protect Mathformer behavior while the framework evolves across Lean, ATP, KB, SymPy, docs, and skills.
%
% Intent:
%   Encode regression checks that catch parser breakage, semantic drift, and missing workflow evidence.
%
% Function:
%   Contains executable tests, fixtures, and manifests used for local and CI validation.
%
% Maintenance Notes:
%   Keep this header synchronized with major responsibility changes, and prefer updating tests or examples when behavior changes.

% mathformer-case-id: wiki_set_inter_union_distrib
% expected-lean-name: WikiLevelGold.wiki_set_inter_union_distrib
% controlled-suite: wiki_level_latex_to_lean
% difficulty: easy
\documentclass[11pt]{article}
\usepackage[T1]{fontenc}
\usepackage[utf8]{inputenc}
\usepackage{lmodern}
\usepackage{amsmath,amssymb,amsthm}

\newtheorem{theorem}{Theorem}

\title{Intersection Distributes over Union}
\author{}
\date{}

\begin{document}

\maketitle

\section*{Controlled Origin Annotation}
This is a synthetic controlled case written in a Wikipedia-style exposition.
It is intended to test extensional set reasoning, membership expansion, and
the translation of informal logical cases into Lean proof branches.

\section*{Statement and Proof}

\begin{theorem}[Intersection distributes over union]
Let $A$, $B$, and $C$ be subsets of a set $X$. Then
\[
  A \cap (B \cup C) = (A \cap B) \cup (A \cap C).
\]
\end{theorem}

\begin{proof}
We prove equality by showing that the two sets have the same elements.
Let $x \in A \cap (B \cup C)$. Then $x \in A$ and $x \in B \cup C$.
If $x \in B$, then $x \in A \cap B$, so
$x \in (A \cap B) \cup (A \cap C)$. If $x \in C$, then
$x \in A \cap C$, so again $x \in (A \cap B) \cup (A \cap C)$.

Conversely, let $x \in (A \cap B) \cup (A \cap C)$. If
$x \in A \cap B$, then $x \in A$ and $x \in B$, hence
$x \in A \cap (B \cup C)$. If $x \in A \cap C$, then $x \in A$ and
$x \in C$, hence $x \in A \cap (B \cup C)$. Thus the two sets are equal.
\end{proof}

\section*{Lean Target}
Formalize the theorem for arbitrary type $\alpha$ and sets
$A,B,C : \operatorname{Set}\,\alpha$. The expected Lean theorem is
\texttt{WikiLevelGold.wiki\_set\_inter\_union\_distrib}.

\end{document}
